\documentclass[a4paper, hidelinks, 12pt]{article}
\input{makra.tex}

\let\leq=\leqslant % redefine
\let\geq=\geqslant % redefine
\renewcommand\qedsymbol{$\blacksquare$}
\newcommand{\contr}{\hfill\Lightning}

\usepackage{ulem}


\begin{document}

\begin{Large}
    \textbf{pmat}
    \hfill
    3 zadanie domowe
    \hfill
    17.10.2024
\end{Large}

\rule{\textwidth}{1pt}


\textbf{Zadanie 5D.}

\medskip 

\textbf{Zadanie.} Prostokąt podzielono odcinkami równoległymi do boków na małe prostokąty.
Następnie wyróżniono każdy punkt, który jest wierzchołkiem dokładnie dwóch małych
prostokątów. Wykaż, że liczba wyróżnionych punktów jest parzysta.

\medskip

\noindent \textbf{Rozwiązanie błędne.} Przeprowadzimy dowód indukcyjny ze względu na liczbę $n$ małych prostokątów, na które podzielony jest duży prostokąt. 

\medskip

Dla $n = 1$ liczba punktów wyróżnionych jest równa zero, więc jest parzysta, co stanowi
bazę indukcji.

\medskip

W kroku indukcyjnym zakładamy, że dla pewnego naturalnego $k \geq 1$ dowolny podział
prostokąta na $k$ małych prostokątów ma parzystą liczbę punktów wyróżnionych i chcemy
udowodnić, implikację: skoro dla że własność tę ma także dowolny podział na $k + 1$ małych prostokątów.

\medskip

Rozważmy dowolny podział prostokąta na $k+1$  małych prostokątów.
Usuńmy jeden z odcinków podziału tak, aby otrzymać pewien podział dużego prostokąta
na $k$ małych prostokątów. Mamy wówczas trzy możliwości dotyczące łącznej liczby
wyróżnionych punktów:
\begin{itemize}
    \item zmniejszyła się o 2 (drugi rysunek) — jeżeli oba końce usuniętego odcinka były
    wyróżnione;
    \item pozostała niezmieniona (trzeci rysunek) — jeżeli jeden z końców usuniętego odcinka był wyróżniony, a drugi nie;
    \item zwiększyła się o 2 (czwarty rysunek) — jeżeli żaden z końców usuniętego odcinka
    nie był wyróżniony.
\end{itemize}

\noindent We wszystkich przypadkach parzystość liczby wyróżnionych punktów nie uległa zmianie (jeżeli przed usunięciem odcinka liczba ta była parzysta, to po jego usunięciu
także). Z założenia indukcyjnego wynika, że przed usunięciem odcinka liczba wyróżnionych punktów była parzysta. To kończy krok indukcyjny.


\medskip

Na mocy zasady indukcji stwierdzamy, że teza zadania została udowodniona dla dowolnego podziału prostokąta.

\textcolor{red}{[Przede wszystkim, gdy bazą indukcji jest najmniejsze możliwe $n$, należy udowodnić $n \rightarrow n+1$, a nie $n+1 \rightarrow n$, ponieważ nie jest to poprawna indukcja. Oprócz tego należy przeprowadzać indukcję po n, a nie po dowolnie wybranym k. Co więcej, w akapicie trzecim zapowiedziane jest wynikanie $n \rightarrow n+1$, które później nie zostaje nigdzie wykonane. Poza tym przeprowadzone rozumowanie jest błędne, nawet gdyby indukcja była poprawna.]}

\textbf{Rozwiązanie poprawne} (mam nadzieję) Dorysowanie odcinka zawsze dwukrotnie zmieni parzystość (czyli nie zmieni wcale). Jeśli koniec odcinka wyląduje na innym odcinku, ten punkt będzie wierzchołkiem dokładnie dwóch utworzonych właśnie prostokątów, zatem dodajemy do wyniku 1. Jeśli jego koniec wyląduje  w miejscu, gdzie koniec innego odcinka spotyka się z odcinkiem, w tamtym punkcie znajdą się teraz wierzchołki 4 prostokątów i dwa końce odcinków, z których żaden nie jest oznaczony. Wcześniej jeden z tych końców był oznaczony, zatem od wyniku odjęliśmy 1. Gdy koniec odcinka wyląduje w punkcie gdzie spotykają się 3 inne końce odcinków, wszystkie 3 z nich są oznaczone. Dodanie tam kolejnego sprawi, że punkt ten stanie się wierzchołkiem 4 prostokątów a nie dwóch, zatem musimy odznaczyć tamten wierzchołek, czyli odjąć 1 od wyniku. W takim razie niezależnie od układu, każdy koniec nowo dodanego wierzchołka zmieni parzystość o 1, zatem cały wierzchołek nie zmieni parzystości. Faktycznie, bazą indukcji może być niepodzielony prostokąt, gdzie wyróżnionych wierzchołków jest 0, jednak przeprowadzamy indukcję po liczbie odcinków a nie małych prostokątów.  


\textbf{Zadanie 6D.} Określić taką rodzinę $\{A_{i,j} \mid i, j \in I\}$, żeby wszystkie poniższe
zbiory były różne:

\[
\bigcup_{i} \bigcap_{j} A_{i,j}, \quad \bigcap_{i} \bigcup_{j} A_{i,j}, \quad \bigcup_{i} \bigcup_{j} A_{i,j}, \quad \bigcap_{i} \bigcap_{j} A_{i,j}, \quad \bigcup_{j} \bigcap_{i} A_{i,j}, \quad \bigcap_{j} \bigcup_{i} A_{i,j}.
\]


\textbf{Rozwiązanie.} Rozważmy liczby $a,b,c,d,e \in \mathbb{R}$, takie że $a\neq b \neq c \neq d \neq e $ oraz rodzinę $\mathcal{A} = A_1,1, A_1,2, A_2,1, A_2,2$ taką że: 
\begin{itemize}
    \item $A_1,1 = \{a, c, d\}$
    \item $A_1,2 = \{b, d\}$
    \item $A_2,1 = \{a, e\}$
    \item $A_2,2 = \{b, c, e\}$
\end{itemize}

Wtedy:
\begin{itemize}
    \item $\bigcup_{i} \bigcap_{j} A_{i,j} = \{d\} \cup \{e\} = \{d, e\}$
    \item $\bigcap_{i} \bigcup_{j} A_{i,j} = \{a, b, c, d\} \cap \{b, c, d, e\} = \{b, c, d\}$
    \item $\bigcup_{i} \bigcup_{j} A_{i,j} = \{a, c, d, e\} \cup \{b, c, d, e\} = \{a, b, c, d, e\}$
    \item $\bigcap_{i} \bigcap_{j} A_{i,j} = \{a\} \cap \{b\} = \emptyset$
    \item $\bigcup_{j} \bigcap_{i} A_{i,j} = \{a\} \cup \{b\} = \{a, b\} $
    \item $\bigcap_{j} \bigcup_{i} A_{i,j} = \{a, c, d, e\} \cap \{b, c, d, e\} = \{c, d, e\}$
\end{itemize}






\textit{Autor: Dominika Piętka}\\

\end{document}